\section{Introduction}
\label{sec:introduction}

This paper develops the Bell-theorem stance of a deterministic brane-lattice substrate --- a 4D
spring-lattice model whose elements are summarized in \Cref{sec:substrate} and constructed in full
in a companion paper \citep{MolzbergerSubstrate}. It gives the constructive case for a retrocausal
worldtube reading of entanglement and traces its consequences for how the model's simulation
program should interpret its own results.

\paragraph{Why Bell forces an interpretive stance.}
The substrate is deterministic (\Cref{eq:brane-eom}) and local (all forces are nearest-neighbor
spring contacts at finite propagation speed).
Bell's theorem then forces an interpretive stance before any specific construction is developed: it
rules out a large class of classical-substrate completions and constrains which interpretations
are available.
The retrocausal worldtube reading is therefore not a stylistic preference but a
\emph{forced consequence} of the substrate's own ontological commitments.

\paragraph{The substrate's stance.}
Bell's theorem rules out theories that are simultaneously (D) deterministic, (MI)
measurement-independent, and (L) local in the sense of finite-speed influence propagation
\citep{Bell1964,CHSH1969}.
The substrate satisfies (D) by construction and (L) by its nearest-neighbor dynamics.
This means (MI) must be relaxed.
Two ways to relax (MI) are available in principle:
\begin{itemize}
  \item \emph{Superdeterminism:} correlate detector settings with the hidden state at some
  common past boundary condition.
  This is logically consistent but explanatorily circular: it assumes the very correlations it
  is supposed to explain \citep{Hossenfelder2020Superdeterminism}.
  \item \emph{Retrocausality:} allow the future measurement context to propagate \emph{backward}
  along the worldtube to the shared past of the entangled pair.
\end{itemize}
The substrate's own ontology closes the superdeterminism route.
Superdeterminism requires a cosmic conspiracy between detector choices and the initial substrate
configuration; but the substrate is a local elastic medium, not a pre-programmed Oracle.
It has no mechanism to coordinate microscopic initial conditions at spacelike separation with
macroscopic measurement choices made later.
The retrocausal reading, by contrast, requires only what the 4D brane ontology already provides:
a 4D-in-4D world-volume in which future boundary conditions are as natural as past ones.
Under (D) + (L) + the rejection of superdeterminism, retrocausality is the completion the
substrate's own dynamics supply --- and the case for it is \emph{constructive}: \Cref{sec:worldtube}
exhibits the mechanism and shows that it reproduces the measured entangled-photon statistics,
rather than resting on having excluded every conceivable alternative.

\paragraph{Paper organization.}
\Cref{sec:substrate} states the substrate model used here.
\Cref{sec:bell-constraint} develops the Bell constraint and the substrate's stance.
\Cref{sec:worldtube} develops the retrocausal worldtube interpretation: time-symmetry of the brane
action, the thermodynamic arrow of time and soliton chirality (matter vs antimatter as
opposite-chirality worldtubes), entanglement as V-branching, the measurement mechanism that
reproduces the entangled-photon correlations (no-signalling, the Tsirelson-bounded $\cos 2(a-b)$
law, and the Born weight as typicality), and implications for the simulation program.